% ---------
% --------
%!TEX TS-program = pdflatex
%!TEX encoding = UTF-8 Unicode
%!TEX spellcheck = English (Aspell)
\documentclass[11pt]{article}
\usepackage{lucy,handout,graphicx}
\usepackage[utf8]{inputenc}		% Allow some non-ASCII Unicode in
\usepackage{tikz}

\def\Cdot{\mathbin{\text{\normalfont \textbullet}}}
\def\Sym#1{\texttt{\color{BrickRed}#1}}
\def\BlueSym#1{\textbf{\texttt{\color{RoyalBlue}#1}}}
\allowdisplaybreaks

\newtheorem{theorem}{Claim}

% =========================================================
%   Define common stuff for solution headers
% =========================================================
\Class{CSCI 332}
\Semester{Fall 2025}

% =========================================================
\begin{document}

% ---------------------------------------------------------


\headers{CSCI 332}{Homework 3 Solutions}{Fall 2026}

\begin{enumerate}

\item (3 points) Consider the number of the syllables in the pronunciation of the integer $n$.
For example, when $n=1$, the pronunciation is ``one'' and there is one syllable.
For each of the following, give the tightest/simplest function possible, or state that
there is no function possible.
\begin{enumerate}
    \item Give a function $f(n)$ so that the number of syllables in the pronunciation of $n$ is $O(f(n))$ for all $n \geq 1$.
    \begin{Solution}
        For large $n$, integers in English are pronounced in groups of three
        digits (hundreds, thousands, millions, billions, etc.).  For example,
        the number $999,999,999$ is pronounced as ``nine hundred ninety-nine
        million, nine hundred ninety-nine thousand, nine hundred ninety-nine''.
        Each group of three digits has a constant maximum number of sylalbles,
        and the number of groups of three digits is proportional to $\log_{10} n$.
        So,
        the tightest upper bound on the number of syllables in the pronunciation
        of $n$ is $O(\log n)$. (So $f(n) = \log n$.)
    
    \end{Solution}
    \item Give a function $g(n)$ so that the number of syllables in the pronunciation of $n$ is $\Omega(g(n))$ for all $n \geq 1$.
    \begin{Solution}
        Even as $n$ grows large, each whole multiple of $1000$ has a short name
        (million, billion, trillion, etc.). Therefore, the tightest lower bound
        on the number of syllables in the pronunciation of $n$ is $\Omega(1)$.
        (So $g(n) = 1$.)
    
    \end{Solution}
    \item Give a function $h(n)$ so that the number of syllables in the pronunciation of $n$ is $\Theta(h(n))$ for all $n \geq 1$.
    \begin{Solution}
        Since the tightest upper bound is $O(\log n)$ and the tightest lower
        bound is $\Omega(1)$, there is no function $h(n)$ so that the number of
        syllables in the pronunciation of $n$ is $\Theta(h(n))$ for all $n \geq
        1$.
    \end{Solution}
\end{enumerate}

\item (3 points) Prove that $\log_a n = \Theta(\log_b n)$ for any positive
constants $a \neq 1, b$ by proving each of the following. You must use the
definition of big-O and big-Omega in your proofs; that is, you must give a $c$
and an $n_0$. (Hint: Use the change of base formula for logarithms, which states
that $\log_a n = \frac{\log_b n}{\log_b a}$.)

    \begin{enumerate}[(a)]
        \item $log_a n = O(log_b n)$
    \begin{Solution}
        To show that $\log_a n = O(\log_b n)$, we need to find constants $c > 0$
        and $n_0 \geq 1$ such that for all $n \geq n_0$, the following
        inequality holds:

        \[
            \log_a n \leq c \cdot \log_b n.
        \]

        Let $c = \frac{1}{\log_b a}$. Since $a > 0$ and $a \neq 1$, we have
        $\log_b a > 0$, so $c > 0$.

        Now notice that

        \begin{align*}
            \log_a n & = \frac{\log_b n}{\log_b a} & \text{(change of base formula)} \\
            & = \frac{1}{\log_b a} \cdot \log_b n & \text{(rewriting)} \\
            & = c \cdot \log_b n & \text{(substituting $c$)}.
        \end{align*}
        
        So it must be the case that $\log_a n \leq c \cdot \log_b n$ for all $n \geq 1$,
        so we can choose $n_0 = 1$ (or any positive integer), and the inequality
        holds for all $n \geq n_0$. Therefore, $\log_a n =
        O(\log_b n)$.

    \end{Solution}
        \item $log_a n = \Omega(log_b n)$
    \begin{Solution}
        In (a), we showed that choosing $c = \frac{1}{\log_b a}$ implies that
        $\log_a n = c \cdot \log_b n$ for all $n \geq 1$. So we can choose the same $c$ and still have
        $\log_a n \geq c \cdot \log_b n$ for all $n \geq 1$, and so we have given a valid choice of $c$ and $n_0$
        to show that $\log_a n = \Omega(\log_b n)$.
    \end{Solution}
    \end{enumerate} 

\item (4 points) For each of the following algorithms, you will analyze both the worst and best case runtimes.
\begin{enumerate}
    \item
    \begin{algo}
        \underline{IsPrime(n):}\\
        \> if $n < 2$\\
        \> \> return False\\
        \> for $d = 2$ to $\lfloor \sqrt{n} \rfloor$\\
        \> \> if $n \bmod d = 0$\\
        \> \> \> return False\\
        \> return True
    \end{algo}
    \begin{enumerate}
        \item Describe a best-case input to IsPrime.
    \begin{Solution}
        We need to identify the class of inputs for which the algorithm runs in the least amount of time.
        Notice that for any even $n$, the algorithm will return false after a constant number of steps.
    \end{Solution}
        \item Give a function $f(n)$ so that the best-case runtime of IsPrime is $\Theta(f(n))$.
    \begin{Solution}
        $f(n) = 1$.
    
    \end{Solution}
        \item Describe a worst-case input to IsPrime.
    \begin{Solution}
        A worst-case input for IsPrime is any prime number, since the algorithm
        will have to check all possible divisors up to $\lfloor \sqrt{n}
        \rfloor$ before returning True.
    
    \end{Solution}
        \item Give a function $f(n)$ so that the worst-case runtime of IsPrime is $\Theta(f(n))$.
    \begin{Solution}
        $f(n) = \sqrt{n}$.
    \end{Solution}
    \end{enumerate}
    \item 
    \begin{algo}
        \underline{HasDuplicate(A[1..n]):}\\
        \> for $i = 1$ to $n$\\
        \> \> for $j = i+1$ to $n$\\
        \> \> \> if $A[i] = A[j]$\\
        \> \> \> \> return True\\
        \> return False
    \end{algo}
    \begin{enumerate}
        \item Describe a best-case input to HasDuplicate.
        \begin{Solution}
            A best-case input for HasDuplicate is an array where the first two
            elements are equal, since the algorithm will return True after a
            constant number of steps.
        \end{Solution}
        \item Give a function $f(n)$ so that the best-case runtime of HasDuplicate is $\Theta(f(n))$.
        \begin{Solution}
            $f(n) = 1$.
        \end{Solution}
        \item Describe a worst-case input to HasDuplicate.
        \begin{Solution}
            A worst-case input for HasDuplicate is an array where all elements
            are distinct, since the algorithm will have to check all pairs of
            elements before returning False.
        \end{Solution}
        \item Give a function $f(n)$ so that the worst-case runtime of HasDuplicate is $\Theta(f(n))$.
        \begin{Solution}
            $f(n) = n^2$.
        \end{Solution}
    \end{enumerate}
\end{enumerate}
    

\end{enumerate}

\end{document}

